\section{What remains open and which reductions are promising}
\label{sec:strategy}

The variable objective, the degree objective, and the coefficient objective
favor different presentations. The 43-variable specification wins the
first comparison among the constructions supplied here. The 45-variable
semilocal version has a much smaller degree bound. The quartic problem
has a separate concrete answer with 1367 variables, or 1391 when a full
ordinary integer coefficient table of moderate bit length is preferred.
These are upper bounds, with explicit constructions and qualified
arithmetic dependencies; none is a proof of optimality.

\subsection{The arithmetic review which still matters}

The most important independent review remains the manuscript's restricted
$2$-converse and its fixed-field coefficient transfer, followed by the
full five-contact height theorem with the stated constant. The new
shared-point argument additionally assumes Sun's Theorem 4.1. We checked
its statement, its order-compatibility consequences, and the distinction
between the unconditional Artin-symbol theorem of Frei and Tinkov\'a
and their conditional primitive-root theorem
\cite{FreiTinkova2025}. Sun's analytic argument adapts polynomial-prime
transference \cite{TaoZiegler2018} and includes a further large-prime
sieve. The present work does not claim an independent reproduction of
all those analytic estimates. Its new arithmetic theorems are explicitly
conditional on the cited statement.

The additional input has a precise purpose. The original private
comparison points need compatibility separately at each selected prime.
One shared index $h$ needs, in addition,
\[
             \gcd(d_{p,-},d_{q,+})=4\qquad\text{for all }p,q.
\]
Otherwise an odd common divisor makes the required residues $0$ and $4$
incompatible. Sun's theorem supplies this cross-prime order condition,
together with avoidance of the primes in the fixed auxiliary progression.
The shared-point construction therefore cannot be justified merely by
reusing the original private witnesses. Its completeness proof and its
new external input are essential.

\subsection{Reducing the number of variables further}

The preferred source has 42 retained coordinates: the input, eighteen
abscissa components for nine elliptic points, two Lucas proxies, one
common slope, and twenty structural parameters for ten pole predicates.
The final finite-algebra polynomial adds one coordinate. Its eliminated
point roots, conic roots, and arithmetic expressions no longer contribute
to this number. Eliminating more finite algebraic roots therefore improves
the degree but does not by itself improve the 43-variable count.

A further variable saving must change this retained base. Possible places
to look are a pole predicate with fewer structural parameters, a rigorous
sharing of parameters between several pole predicates, or a height and
Lucas argument using fewer elliptic indices. Such sharing would need a
simultaneous rational-solubility theorem; a dimension count or the local
solubility of each predicate separately is insufficient. The conic
obstruction proved here also rules out a tempting shortcut: a semilocal
integrality set containing all integers cannot be represented by a finite
union of equations quadratic in only two witnesses, even if arbitrary
input degree is allowed.

No construction with ten variables has been proved here, and no lower
bound excluding ten variables has been proved. The achieved 43-variable
construction lies within the requested range of ten to fifty, under the
stated arithmetic inputs. The gap to the unrestricted lower bound of
three total variables remains substantial.

\subsection{Degree and coefficient size at a fixed variable count}

The enormous degree bound for 43 variables is dominated by the fixed
elliptic exclusion product. Its cutoff depends on the explicit height
constant, not on the finite-algebra rank alone. A materially better
height constant, or a less costly zero-witness way of imposing the needed
index size, would have disproportionate value. The 45-variable alternative
avoids that product, which explains its much smaller degree bound.
The conic parametrizations already reduce the largest algebra rank to
$2^{29}$ or $2^{30}$. Further use of sparse norm formulas or more economical
separating-element constructions could improve the present conservative
bounds without changing the number of variables.

For quartics, direct two-coordinate field relations, shared polynomial
coefficients, blocked polynomial evaluation, and quadratic-preserving
elimination have all produced actual savings. Further searches should
operate on the final residual system and retain exact reconstruction
certificates. At a fixed coordinate count, rational diagonal changes of
variables and changes of residual basis may reduce coefficient height.
The comparison must use primitive integer content after clearing
numerical denominators. Introducing extra coordinates for constants is a
different trade: it improves coefficient size but worsens the first
objective, which is why both quartics are supplied.

\subsection{Lower bounds and number fields}

The remaining two-witness case includes sign-changing quartic surfaces.
The degree-three unirationality argument and the conic-specialization
argument do not control all such surfaces; smooth quartic surfaces are
K3 surfaces. The quadratic-system lower bound likewise encounters K3
surfaces at intersections of three quadrics. Any stronger unconditional
variable lower bound must address geometry beyond the cases proved here.

The exact $\Q(i)$ section answers the counterexample search for the
released equation, and the bounded-extension obstruction shows that no
single fixed integer-defining equation can work unchanged in every number
field. Neither proves the existence of a number field whose ring of
integers admits no Diophantine definition in that field. The unresolved
question is field-dependent. Under rational integrality, the equivalence
in Theorem~\ref{thm:nf-equivalence} identifies the missing step with
existential definability of the embedded rational subfield. MRDP and
undecidability for the ring of integers do not supply that step.

\section{A conditional geometric obstruction}
\label{sec:mazur}

The numerical lower bounds and the conic obstruction above are
unconditional. There is also a useful
conjectural obstruction of a different kind. Mazur's finite-component
conjecture states that the closure of the rational points of a variety
over $\Q$ in its real points has finitely many connected components
\cite[Conjecture 3]{Mazur1995}. It would rule out a Diophantine
definition of $\Z$ with any degree or number of variables. It must
not be used as an established theorem against the proposed source.

\begin{proposition}
If a polynomial defines $\Z$ in $\Q$, the real closure of its rational
zero set has infinitely many connected components.
\end{proposition}
\begin{proof}
Let $V$ be its affine zero set and let $C$ be the closure of $V(\Q)$
inside $V(\mathbb R)$. The input projection is continuous. Its value
on every convergent sequence of rational points is a limit of integers,
and hence an integer; therefore its image on $C$ is contained in $\Z$.
Every integer occurs already on $V(\Q)$. A connected component of
$C$ has connected image in the discrete set $\Z$, hence has only one
input value. Infinitely many components are necessary.
\end{proof}

Mazur's formulation uses reduced schemes of finite type over $\Q$,
so it applies directly to the reduced affine zero set used here
\cite[Section 1 and Conjecture 3]{Mazur1995}. This topological obstruction
clarifies the strength of the proposed theorem. It does not add a
premise to our exact algebraic transformations.
